\documentclass[11pt]{article}

\usepackage[margin=1in]{geometry}
\usepackage{amsmath, amssymb, mathtools}
\usepackage{xcolor}
\usepackage{enumitem}
\usepackage{array}
\usepackage{hyperref}

\hypersetup{
  colorlinks=true,
  linkcolor=blue,
  urlcolor=blue,
  citecolor=blue
}

% -----------------------------------------------------------------------------
% Basic notation
% -----------------------------------------------------------------------------

\newcommand{\R}{\mathbb{R}}
\newcommand{\E}{\mathbb{E}}
\newcommand{\Prob}{\mathbb{P}}
\newcommand{\one}{\mathbf{1}}

\newcommand{\Dir}{\operatorname{Dirichlet}}
\newcommand{\Cat}{\operatorname{Categorical}}
\newcommand{\softmax}{\operatorname{softmax}}
\newcommand{\softplus}{\operatorname{softplus}}
\newcommand{\argmax}{\operatorname*{arg\,max}}
\newcommand{\mean}{\operatorname{mean}}
\newcommand{\stopgrad}{\operatorname{stopgrad}}

\newcommand{\Tree}{\mathcal{T}}
\newcommand{\Path}{\mathcal{P}}
\newcommand{\States}{\mathcal{S}}
\newcommand{\Actions}{\mathcal{A}}
\newcommand{\Legal}{\mathcal{A}_{\mathrm{legal}}}
\newcommand{\Z}{\mathcal{Z}}

\newcommand{\Step}{\operatorname{Step}}
\newcommand{\Eval}{\operatorname{Eval}}
\newcommand{\Flip}{\operatorname{flip}}
\newcommand{\Align}{\operatorname{align}}
\newcommand{\Utility}{U}

\newcommand{\ThompsonSelect}{\operatorname{ThompsonSelect}}
\newcommand{\EdgeBase}{\operatorname{EdgeBase}}
\newcommand{\EdgePosterior}{\operatorname{EdgePosterior}}
\newcommand{\StateSearchPosterior}{\operatorname{StateSearchPosterior}}
\newcommand{\BackupPath}{\operatorname{BackupPath}}
\newcommand{\NextRequest}{\operatorname{NextRequest}}
\newcommand{\ConsumeResult}{\operatorname{ConsumeResult}}
\newcommand{\BuildEvalBatch}{\operatorname{BuildEvalBatch}}
\newcommand{\GpuEvaluate}{\operatorname{GpuEvaluate}}
\newcommand{\PosteriorBestPolicyTarget}{\operatorname{PosteriorBestPolicyTarget}}
\newcommand{\FinishSearch}{\operatorname{FinishSearch}}
\newcommand{\MakeTargets}{\operatorname{MakeTargets}}

% -----------------------------------------------------------------------------
% Lightweight algorithm box
% -----------------------------------------------------------------------------

\newenvironment{algorithmblock}[1]
{
  \par\medskip
  \noindent\hrule\par\smallskip
  \noindent\textbf{#1}
  \par\smallskip
}
{
  \par\smallskip\noindent\hrule\par\medskip
}

\newcommand{\Input}{\textbf{Input: }}
\newcommand{\Output}{\textbf{Output: }}
\newcommand{\Require}{\textbf{Require: }}
\newcommand{\Return}{\textbf{return }}
\newcommand{\If}{\textbf{if }}
\newcommand{\Then}{\textbf{ then }}
\newcommand{\Else}{\textbf{else }}
\newcommand{\EndIf}{\textbf{end if}}
\newcommand{\For}{\textbf{for }}
\newcommand{\EndFor}{\textbf{end for}}
\newcommand{\While}{\textbf{while }}
\newcommand{\EndWhile}{\textbf{end while}}
\newcommand{\Do}{\textbf{ do }}
\newcommand{\Break}{\textbf{break}}
\newcommand{\Continue}{\textbf{continue}}

\title{Dirichlet-Q AlphaZero: Posterior Tree Search Algorithms}
\author{}
\date{}

\begin{document}

\maketitle

\section*{1. Purpose of this file}

This document defines the search-side algorithms for Dirichlet-Q AlphaZero.

The central change from standard AlphaZero-style MCTS is that the tree does not
store scalar action values as its main memory. Instead, each tree edge stores a
WDL Dirichlet posterior. Search is therefore a posterior-refinement process.

The CPU owns a mutable posterior tree. The GPU only evaluates batches of leaf
states. The search interface is:

\[
\NextRequest(\Tree),
\qquad
\ConsumeResult(\Tree,r),
\qquad
\FinishSearch(\Tree).
\]

The CPU repeatedly selects paths through the current posterior tree, creates
leaf-evaluation requests, sends leaf states to the GPU in batches, and then backs
up WDL evidence when GPU results return.

This version deliberately does \emph{not} include temporary reservation mass or
virtual-loss-style in-flight Dirichlet mass. In-flight leaves may still be tracked
for scheduling, but they do not alter any edge posterior.

\section*{2. Outcome space and utility}

For a two-player zero-sum game, define the WDL outcome space

\[
\Z = \{L,D,W\}.
\]

A WDL distribution is

\[
\phi = (\phi_L,\phi_D,\phi_W) \in \Delta^2.
\]

The scalar utility used everywhere is

\[
\Utility(\phi) = \phi_W - \phi_L.
\]

For a Dirichlet parameter vector

\[
\alpha = (\alpha_L,\alpha_D,\alpha_W),
\]

define its mean as

\[
\mean(\alpha)
=
\frac{\alpha}{\alpha_0},
\qquad
\alpha_0 = \sum_{z \in \{L,D,W\}} \alpha_z.
\]

When the intended action rule is greedy with respect to a posterior mean, define

\[
\bar\phi_{v,a}
=
\mean
\left(
\alpha^{\mathrm{post}}_{v,a}
\right).
\]

Then the greedy mean action is

\[
a^{\star}_{\mathrm{mean}}
=
\argmax_{a\in\Legal(s_v)}
\left(
\bar\phi_{v,a,W}-\bar\phi_{v,a,L}
\right).
\]

This is different from Thompson selection, where the maximization is over a
fresh posterior sample rather than over the posterior mean.

The opponent-perspective flip is

\[
\Flip(d_L,d_D,d_W)
=
(d_W,d_D,d_L).
\]

The same flip applies to Dirichlet parameters:

\[
\Flip(\alpha_L,\alpha_D,\alpha_W)
=
(\alpha_W,\alpha_D,\alpha_L).
\]

For zero-sum alternating-turn games,

\[
\Utility(\Flip(d)) = -\Utility(d).
\]

\section*{3. Network outputs}

At state \(s\), the network predicts:

\[
\ell_\theta(s,a),
\qquad
\alpha_\theta^V(s),
\qquad
\alpha_\theta^Q(s,a).
\]

The policy head is

\[
\pi_\theta(a \mid s)
=
\softmax(\ell_\theta(s,\cdot))_a.
\]

The value head represents a WDL belief over the state value:

\[
\phi_s^V
\sim
\Dir(\alpha_\theta^V(s)).
\]

The Q head represents a WDL belief over each state-action value:

\[
\phi_{s,a}^Q
\sim
\Dir(\alpha_\theta^Q(s,a)).
\]

Both Dirichlet heads use the stable mean-concentration parameterization:

\[
\alpha_\theta
=
\alpha_{\theta,0}\bar{\phi}_\theta,
\]

where

\[
\bar{\phi}_\theta = \softmax(r_\theta),
\qquad
\alpha_{\theta,0} = \softplus(t_\theta)^2.
\]

\section*{4. Posterior tree state}

A posterior tree is denoted by

\[
\Tree.
\]

Each node \(v\) stores:

\[
s_v,
\qquad
\texttt{expanded}(v),
\qquad
\alpha_\theta^V(s_v),
\qquad
\alpha_\theta^Q(s_v,\cdot),
\qquad
\ell_\theta(s_v,\cdot).
\]

Each edge \((v,a)\) stores:

\[
\texttt{child}(v,a),
\qquad
E_{v,a} \in \R_{\ge 0}^{3},
\qquad
N_{v,a} \in \mathbb{N}.
\]

Here

\[
E_{v,a}
=
\text{completed WDL pseudo-evidence accumulated on edge } (v,a).
\]

The posterior used for both selection and targets is

\[
\alpha^{\mathrm{post}}_{v,a}
=
\alpha^{\mathrm{base}}_{v,a}
+
E_{v,a}.
\]

No separate selection posterior is used. In-flight scheduling state may exist in
an implementation, but it is not part of the Bayesian tree state and it is not
added to any Dirichlet posterior.

\section*{5. Edge base prior}

For an edge \((v,a)\), the base Dirichlet prior depends on whether the child
state has already been expanded.

If the child is not expanded, use the Q-head prediction at the parent node:

\[
\alpha^{\mathrm{base}}_{v,a}
=
\alpha_\theta^Q(s_v,a).
\]

If the child \(u=\texttt{child}(v,a)\) is expanded, use the child value head,
aligned back to the perspective of the player to move at node \(v\):

\[
\alpha^{\mathrm{base}}_{v,a}
=
\Align_{v \leftarrow u}
\left(
\alpha_\theta^V(s_u)
\right).
\]

For ordinary alternating-turn two-player games, this is usually

\[
\alpha^{\mathrm{base}}_{v,a}
=
\Flip
\left(
\alpha_\theta^V(s_u)
\right).
\]

The general alignment notation is used because some environments may include
passes, chance transitions, or non-alternating turns.

\begin{algorithmblock}{Algorithm 1: EdgeBase\((\Tree,v,a)\)}

\Input posterior tree \(\Tree\), expanded node \(v\), legal action \(a\)

\begin{enumerate}[leftmargin=2em]
  \item If edge \((v,a)\) has an expanded child \(u\), set
  \[
  \alpha^{\mathrm{base}}_{v,a}
  \gets
  \Align_{v \leftarrow u}
  \left(
  \alpha_\theta^V(s_u)
  \right).
  \]

  \item Otherwise set
  \[
  \alpha^{\mathrm{base}}_{v,a}
  \gets
  \alpha_\theta^Q(s_v,a).
  \]

  \item \Return \(\alpha^{\mathrm{base}}_{v,a}\).
\end{enumerate}

\end{algorithmblock}

\begin{algorithmblock}{Algorithm 2: EdgePosterior\((\Tree,v,a)\)}

\Input posterior tree \(\Tree\), expanded node \(v\), legal action \(a\)

\begin{enumerate}[leftmargin=2em]
  \item Compute
  \[
  \alpha^{\mathrm{base}}_{v,a}
  \gets
  \EdgeBase(\Tree,v,a).
  \]

  \item Compute
  \[
  \alpha^{\mathrm{post}}_{v,a}
  \gets
  \alpha^{\mathrm{base}}_{v,a}
  +
  E_{v,a}.
  \]

  \item \Return \(\alpha^{\mathrm{post}}_{v,a}\).
\end{enumerate}

\end{algorithmblock}

\section*{6. Thompson action selection at a node}

At an expanded node \(v\), selection samples one WDL distribution per legal
action:

\[
\phi_{v,a}
\sim
\Dir
\left(
\alpha^{\mathrm{post}}_{v,a}
\right),
\qquad
 a \in \Legal(s_v).
\]

The selected action is

\[
a^\star
=
\argmax_{a \in \Legal(s_v)}
\Utility(\phi_{v,a})
=
\argmax_{a \in \Legal(s_v)}
\left(
\phi_{v,a,W}-\phi_{v,a,L}
\right).
\]

This is Thompson sampling over WDL action-value beliefs. It is not the same as
choosing the action with largest posterior mean utility, which would use
\(\bar\phi_{v,a}=\mean(\alpha^{\mathrm{post}}_{v,a})\).

\begin{algorithmblock}{Algorithm 3: ThompsonSelect\((\Tree,v)\)}

\Input posterior tree \(\Tree\), expanded node \(v\)

\begin{enumerate}[leftmargin=2em]
  \item For each legal action \(a \in \Legal(s_v)\), compute
  \[
  \alpha_{v,a}
  \gets
  \EdgePosterior(\Tree,v,a).
  \]

  \item Sample
  \[
  \phi_{v,a}
  \sim
  \Dir(\alpha_{v,a})
  \qquad
  \text{for each } a \in \Legal(s_v).
  \]

  \item Compute
  \[
  u_a
  \gets
  \Utility(\phi_{v,a})
  =
  \phi_{v,a,W}-\phi_{v,a,L}
  \qquad
  \text{for each } a \in \Legal(s_v).
  \]

  \item Select
  \[
  a^\star
  \gets
  \argmax_{a \in \Legal(s_v)} u_a.
  \]

  \item \Return \(a^\star\).
\end{enumerate}

\end{algorithmblock}

\section*{7. Leaf evaluation}

Search stops at either terminal or non-terminal leaves.

If the leaf is terminal with outcome \(z\), the evidence is one-hot:

\[
d_\ell = e_z,
\qquad
\lambda_\ell = c_{\mathrm{terminal}}.
\]

If the leaf is non-terminal, the GPU evaluates the leaf state and returns
\(\alpha_\theta^V(s_\ell)\). The WDL evidence is the value-head mean:

\[
d_\ell
=
\mean
\left(
\alpha_\theta^V(s_\ell)
\right),
\qquad
\lambda_\ell = c_{\mathrm{leaf}}.
\]

Usually,

\[
c_{\mathrm{terminal}} > c_{\mathrm{leaf}}.
\]

Terminal evidence is more reliable than bootstrapped neural evidence.

\section*{8. Search-weighted state posterior for backup}

A simple path backup propagates only the leaf evidence \(\lambda_\ell d_\ell\)
to every traversed edge. That is easy, but it ignores the search information
already accumulated below intermediate nodes on the path.

For a more precise backup, define the search-improved state posterior at an
expanded non-terminal node \(u\) by weighting outgoing action posteriors by the
local posterior-best policy:

\[
\beta^V_u
=
\sum_{b\in\Legal(s_u)}
\hat\pi_{\mathrm{search}}^u(b)
\alpha^{\mathrm{post}}_{u,b}.
\]

Equivalently,

\[
\beta^V_u
=
\sum_{b\in\Legal(s_u)}
\hat\pi_{\mathrm{search}}^u(b)
\left(
\alpha^{\mathrm{base}}_{u,b}+E_{u,b}
\right).
\]

This is the precise form of the idea ``propagate \(E+\alpha\), weighted by
\(\pi_{\mathrm{search}}\).'' The object \(\beta^V_u\) is a Dirichlet-parameter
proxy for the value of state \(u\) under its search-improved policy.

The policy \(\hat\pi_{\mathrm{search}}^u\) is computed from the outgoing action
posteriors:

\[
\hat\pi_{\mathrm{search}}^u
=
\PosteriorBestPolicyTarget
\left(
\left\{\alpha^{\mathrm{post}}_{u,b}\right\}_{b\in\Legal(s_u)},
M_{\mathrm{backup}},
\Legal(s_u)
\right).
\]

To add this state posterior to an incoming edge, convert it into pseudo-evidence
with a calibration coefficient \(c_{\mathrm{state}}>0\):

\[
\Delta E_{p,a}^{\mathrm{state}}
=
 c_{\mathrm{state}}
\Align_{p\leftarrow u}
\left(
\beta^V_u
\right).
\]

Using \(c_{\mathrm{state}}=1\) copies the full child state posterior as evidence.
In neural search this may be too confident because subtree evaluations are
correlated and bootstrapped. A safer first setting is to treat
\(c_{\mathrm{state}}\) as a small calibration knob.

\begin{algorithmblock}{Algorithm 4: StateSearchPosterior\((\Tree,u,M_{\mathrm{backup}})\)}

\Input posterior tree \(\Tree\), expanded non-terminal node \(u\), posterior-best sample count \(M_{\mathrm{backup}}\)

\Output search-weighted state posterior \(\beta^V_u\)

\begin{enumerate}[leftmargin=2em]
  \item For each legal action \(b\in\Legal(s_u)\), compute
  \[
  \alpha^{\mathrm{post}}_{u,b}
  \gets
  \EdgePosterior(\Tree,u,b).
  \]

  \item Compute the local posterior-best policy:
  \[
  \hat\pi_{\mathrm{search}}^u
  \gets
  \PosteriorBestPolicyTarget
  \left(
  \left\{\alpha^{\mathrm{post}}_{u,b}\right\}_{b\in\Legal(s_u)},
  M_{\mathrm{backup}},
  \Legal(s_u)
  \right).
  \]

  \item Compute
  \[
  \beta^V_u
  \gets
  \sum_{b\in\Legal(s_u)}
  \hat\pi_{\mathrm{search}}^u(b)
  \alpha^{\mathrm{post}}_{u,b}.
  \]

  \item \Return \(\beta^V_u\).
\end{enumerate}

\end{algorithmblock}

\section*{9. Backup}

Given a path

\[
\Path=((v_0,a_0),\ldots,(v_{\ell-1},a_{\ell-1})),
\]

leaf WDL evidence \(d_\ell\), and evidence strength \(\lambda_\ell\), backup
updates the final edge using the leaf evidence, then propagates search-weighted
state posteriors through the remaining ancestors.

The final edge receives the direct leaf evidence:

\[
E_{v_{\ell-1},a_{\ell-1}}
\gets
E_{v_{\ell-1},a_{\ell-1}}
+
\lambda_\ell
\Align_{v_{\ell-1}\leftarrow \ell}(d_\ell).
\]

For an ancestor edge \((v_i,a_i)\) with expanded child \(u_i=\texttt{child}(v_i,a_i)\),
use the child's search-weighted state posterior:

\[
E_{v_i,a_i}
\gets
E_{v_i,a_i}
+
 c_{\mathrm{state}}
\Align_{v_i\leftarrow u_i}
\left(
\beta^V_{u_i}
\right),
\qquad
\beta^V_{u_i}=\StateSearchPosterior(\Tree,u_i,M_{\mathrm{backup}}).
\]

In ordinary alternating-turn zero-sum games, each alignment is repeated
application of \(\Flip\) depending on the player-to-move perspective.

\begin{algorithmblock}{Algorithm 5: BackupPath\((\Tree,\Path,d_\ell,\lambda_\ell,M_{\mathrm{backup}},c_{\mathrm{state}})\)}

\Input posterior tree \(\Tree\), path \(\Path\), leaf evidence \(d_\ell\), leaf evidence strength \(\lambda_\ell\), posterior-best sample count \(M_{\mathrm{backup}}\), state-backup scale \(c_{\mathrm{state}}\)

\begin{enumerate}[leftmargin=2em]
  \item If \(\Path\) is empty, \Return \(\Tree\).

  \item Let \((v_{\ell-1},a_{\ell-1})\) be the final edge in \(\Path\).

  \item Align the leaf evidence to the final parent:
  \[
  d_{\ell-1}
  \gets
  \Align_{v_{\ell-1}\leftarrow \ell}(d_\ell).
  \]

  \item Add direct leaf evidence to the final edge:
  \[
  E_{v_{\ell-1},a_{\ell-1}}
  \gets
  E_{v_{\ell-1},a_{\ell-1}}
  +
  \lambda_\ell d_{\ell-1}.
  \]

  \item Increment the final edge visit counter:
  \[
  N_{v_{\ell-1},a_{\ell-1}}
  \gets
  N_{v_{\ell-1},a_{\ell-1}}+1.
  \]

  \item For \(i=\ell-2,\ldots,0\):
  \begin{enumerate}[leftmargin=2em]
    \item Let
    \[
    u_i \gets \texttt{child}(v_i,a_i).
    \]

    \item Compute the child search-weighted state posterior:
    \[
    \beta^V_{u_i}
    \gets
    \StateSearchPosterior(\Tree,u_i,M_{\mathrm{backup}}).
    \]

    \item Add calibrated, aligned state posterior evidence:
    \[
    E_{v_i,a_i}
    \gets
    E_{v_i,a_i}
    +
    c_{\mathrm{state}}
    \Align_{v_i\leftarrow u_i}
    \left(
    \beta^V_{u_i}
    \right).
    \]

    \item Increment the edge visit counter:
    \[
    N_{v_i,a_i}
    \gets
    N_{v_i,a_i}+1.
    \]
  \end{enumerate}

  \item \Return \(\Tree\).
\end{enumerate}

\end{algorithmblock}

\section*{10. Initializing a posterior tree}

The root node must be evaluated before search begins, because Thompson selection
requires root priors.

\begin{algorithmblock}{Algorithm 6: InitializePosteriorTree\((s,\theta)\)}

\Input root state \(s\), network parameters \(\theta\)

\begin{enumerate}[leftmargin=2em]
  \item Run the network on the root state:
  \[
  \left(
  \ell_\theta(s,\cdot),
  \alpha_\theta^V(s),
  \alpha_\theta^Q(s,\cdot)
  \right)
  \gets
  f_\theta(s).
  \]

  \item Create root node \(v_{\mathrm{root}}\) with state \(s\).

  \item Store
  \[
  \ell_\theta(s,\cdot),
  \qquad
  \alpha_\theta^V(s),
  \qquad
  \alpha_\theta^Q(s,\cdot)
  \]
  in the root node.

  \item Mark the root node as expanded.

  \item Initialize all edge evidence to zero:
  \[
  E_{v,a} \gets 0.
  \]

  \item Initialize counters:
  \[
  T_{\mathrm{done}}(\Tree) \gets 0,
  \qquad
  T_{\mathrm{inflight}}(\Tree) \gets 0.
  \]

  \item \Return posterior tree \(\Tree\).
\end{enumerate}

\end{algorithmblock}

\section*{11. CPU iterator: producing one leaf-evaluation request}

The CPU tree exposes an iterator-like method:

\[
\NextRequest(\Tree).
\]

It walks down the current posterior tree by Thompson sampling until it finds
either a terminal leaf or a non-terminal leaf that requires GPU evaluation.

If it reaches a terminal leaf, it immediately backs up terminal evidence and
does not emit a GPU request.

If it reaches a non-terminal unexpanded leaf, it returns an evaluation request.
No reservation mass is added to any traversed edge.

An evaluation request contains:

\[
\texttt{tree\_id},
\qquad
\texttt{leaf\_id},
\qquad
\Path,
\qquad
s_\ell.
\]

Here \(\Path\) is the selected path. The leaf node itself can be marked as
in-flight so that the same unevaluated child is not submitted twice.

\begin{algorithmblock}{Algorithm 7: NextRequest\((\Tree,M_{\mathrm{backup}},c_{\mathrm{state}})\)}

\Input posterior tree \(\Tree\), backup posterior-best sample count \(M_{\mathrm{backup}}\), state-backup scale \(c_{\mathrm{state}}\)

\Output either no GPU request, or an evaluation request

\begin{enumerate}[leftmargin=2em]
  \item Set
  \[
  v \gets v_{\mathrm{root}},
  \qquad
  \Path \gets ().
  \]

  \item \While true \Do

  \begin{enumerate}[leftmargin=2em]
    \item \If \(s_v\) is terminal with outcome \(z\) \Then

    \[
    d_\ell \gets e_z,
    \qquad
    \lambda_\ell \gets c_{\mathrm{terminal}}.
    \]

    Run

    \[
    \Tree \gets
    \BackupPath(\Tree,\Path,d_\ell,\lambda_\ell,M_{\mathrm{backup}},c_{\mathrm{state}}).
    \]

    Increment

    \[
    T_{\mathrm{done}}(\Tree)
    \gets
    T_{\mathrm{done}}(\Tree)+1.
    \]

    \Return no GPU request.

    \EndIf

    \item \If node \(v\) is not expanded \Then

    This node is already in flight. The implementation may restart from the
    root, try another active tree, or return no request.

    \Return no GPU request.

    \EndIf

    \item Select an action by Thompson sampling:
    \[
    a
    \gets
    \ThompsonSelect(\Tree,v).
    \]

    \item Append the edge to the path:
    \[
    \Path
    \gets
    \Path \mathbin{\Vert} ((v,a)).
    \]

    \item \If edge \((v,a)\) has an expanded child \(u\) \Then
    \[
    v \gets u.
    \]
    \Continue.
    \EndIf

    \item \If edge \((v,a)\) has a child \(u\) that exists but is not expanded \Then

    The leaf is already in flight. Return no request or restart from the root.
    No posterior mass is changed.

    \Return no GPU request.

    \EndIf

    \item Otherwise create a new child state:
    \[
    s_u \gets \Step(s_v,a).
    \]

    \item Create a child node \(u\) with state \(s_u\), mark it as not expanded,
    and store it as
    \[
    \texttt{child}(v,a)\gets u.
    \]

    \item \If \(s_u\) is terminal with outcome \(z\) \Then

    \[
    d_\ell \gets e_z,
    \qquad
    \lambda_\ell \gets c_{\mathrm{terminal}}.
    \]

    Run

    \[
    \Tree \gets
    \BackupPath(\Tree,\Path,d_\ell,\lambda_\ell,M_{\mathrm{backup}},c_{\mathrm{state}}).
    \]

    Increment

    \[
    T_{\mathrm{done}}(\Tree)
    \gets
    T_{\mathrm{done}}(\Tree)+1.
    \]

    \Return no GPU request.

    \EndIf

    \item Mark \(u\) as in flight and increment
    \[
    T_{\mathrm{inflight}}(\Tree)
    \gets
    T_{\mathrm{inflight}}(\Tree)+1.
    \]

    \item \Return
    \[
    \operatorname{EvalRequest}
    \left(
    \texttt{tree\_id},
    u,
    \Path,
    s_u
    \right).
    \]

  \end{enumerate}

  \item \EndWhile
\end{enumerate}

\end{algorithmblock}

\section*{12. Consuming a GPU result}

A GPU result for leaf \(u\) contains

\[
\ell_\theta(s_u,\cdot),
\qquad
\alpha_\theta^V(s_u),
\qquad
\alpha_\theta^Q(s_u,\cdot).
\]

The CPU stores these outputs in the leaf node, marks the node as expanded, and
backs up neural value evidence.

\begin{algorithmblock}{Algorithm 8: ConsumeResult\((\Tree,r,M_{\mathrm{backup}},c_{\mathrm{state}})\)}

\Input posterior tree \(\Tree\), GPU result \(r\), backup posterior-best sample count \(M_{\mathrm{backup}}\), state-backup scale \(c_{\mathrm{state}}\)

\begin{enumerate}[leftmargin=2em]
  \item Extract from \(r\):
  \[
  u,
  \qquad
  \Path,
  \qquad
  \ell_\theta(s_u,\cdot),
  \qquad
  \alpha_\theta^V(s_u),
  \qquad
  \alpha_\theta^Q(s_u,\cdot).
  \]

  \item Store the network outputs in node \(u\).

  \item Mark node \(u\) as expanded and no longer in flight.

  \item Convert the leaf value Dirichlet into WDL evidence:
  \[
  d_\ell
  \gets
  \mean
  \left(
  \alpha_\theta^V(s_u)
  \right).
  \]

  \item Set neural leaf evidence strength:
  \[
  \lambda_\ell
  \gets
  c_{\mathrm{leaf}}.
  \]

  \item Back up the neural evidence:
  \[
  \Tree
  \gets
  \BackupPath(\Tree,\Path,d_\ell,\lambda_\ell,M_{\mathrm{backup}},c_{\mathrm{state}}).
  \]

  \item Update counters:
  \[
  T_{\mathrm{inflight}}(\Tree)
  \gets
  T_{\mathrm{inflight}}(\Tree)-1,
  \]
  \[
  T_{\mathrm{done}}(\Tree)
  \gets
  T_{\mathrm{done}}(\Tree)+1.
  \]

  \item \Return \(\Tree\).
\end{enumerate}

\end{algorithmblock}

\section*{13. Batched CPU/GPU search driver}

The driver maintains many active posterior trees. This gives game-parallelism.
Each tree may also have several in-flight leaves. This gives leaf-parallelism
without adding reservation mass to the posterior.

Let

\[
G = \text{number of active root trees},
\]

\[
L = \text{maximum in-flight leaves per tree},
\]

\[
F = \text{GPU inference batch size},
\]

and

\[
T = \text{completed simulations per root}.
\]

The available inference batch size is approximately

\[
\min(GL,F).
\]

The GPU evaluator is the compilation boundary. It receives a preallocated batch
of observations and returns a preallocated batch of network outputs.

\begin{algorithmblock}{Algorithm 9: BuildEvalBatch\((\{\Tree_b\}_{b=1}^{G},F,L,M_{\mathrm{backup}},c_{\mathrm{state}})\)}

\Input active trees \(\{\Tree_b\}_{b=1}^{G}\), GPU batch size \(F\), in-flight limit \(L\), backup posterior-best sample count \(M_{\mathrm{backup}}\), state-backup scale \(c_{\mathrm{state}}\)

\Output evaluation request batch \(\mathcal{Q}\)

\begin{enumerate}[leftmargin=2em]
  \item Initialize
  \[
  \mathcal{Q} \gets ().
  \]

  \item \While \(|\mathcal{Q}| < F\) and some tree has remaining budget \Do

  \begin{enumerate}[leftmargin=2em]
    \item Choose an active tree \(\Tree_b\).

    \item \If
    \[
    T_{\mathrm{done}}(\Tree_b) \ge T
    \]
    \Then skip this tree. \EndIf

    \item \If
    \[
    T_{\mathrm{inflight}}(\Tree_b) \ge L
    \]
    \Then skip this tree. \EndIf

    \item Run
    \[
    q
    \gets
    \NextRequest(\Tree_b,M_{\mathrm{backup}},c_{\mathrm{state}}).
    \]

    \item \If \(q\) is a GPU evaluation request \Then append \(q\) to \(\mathcal{Q}\). \EndIf

    \item If no tree can produce more requests, \Break.

  \end{enumerate}

  \item \EndWhile

  \item \Return \(\mathcal{Q}\).
\end{enumerate}

\end{algorithmblock}

\begin{algorithmblock}{Algorithm 10: RunPosteriorTreeSearch\((\{s_b\}_{b=1}^{G},\theta,T,F,L,M_{\mathrm{backup}},c_{\mathrm{state}})\)}

\Input root states \(\{s_b\}_{b=1}^{G}\), network parameters \(\theta\), completed simulations \(T\), GPU batch size \(F\), in-flight limit \(L\), backup posterior-best sample count \(M_{\mathrm{backup}}\), state-backup scale \(c_{\mathrm{state}}\)

\Output completed posterior trees \(\{\Tree_b\}_{b=1}^{G}\)

\begin{enumerate}[leftmargin=2em]
  \item For each root state \(s_b\), initialize
  \[
  \Tree_b
  \gets
  \operatorname{InitializePosteriorTree}(s_b,\theta).
  \]

  \item \While some tree satisfies
  \[
  T_{\mathrm{done}}(\Tree_b) < T
  \]
  \Do

  \begin{enumerate}[leftmargin=2em]
    \item Build a GPU evaluation batch:
    \[
    \mathcal{Q}
    \gets
    \BuildEvalBatch
    \left(
    \{\Tree_b\}_{b=1}^{G},
    F,
    L,
    M_{\mathrm{backup}},
    c_{\mathrm{state}}
    \right).
    \]

    \item \If \(\mathcal{Q}\) is empty \Then \Continue. \EndIf

    \item Evaluate all requested leaves on the GPU:
    \[
    \mathcal{R}
    \gets
    \GpuEvaluate(\mathcal{Q},\theta).
    \]

    \item For each result \(r\in\mathcal{R}\), send it back to its owning CPU tree:
    \[
    \Tree_{r}
    \gets
    \ConsumeResult(\Tree_{r},r,M_{\mathrm{backup}},c_{\mathrm{state}}).
    \]

  \end{enumerate}

  \item \EndWhile

  \item \Return \(\{\Tree_b\}_{b=1}^{G}\).
\end{enumerate}

\end{algorithmblock}

\section*{14. Posterior-best policy target}

After search, the root action posterior is

\[
\alpha_{\mathrm{root}}(a)
=
\EdgePosterior(\Tree,v_{\mathrm{root}},a).
\]

The search-improved policy target is the posterior probability that each action
is optimal:

\[
\pi_{\mathrm{search}}(a\mid s)
=
\Prob
\left(
a
=
\argmax_{b\in\Legal(s)}
\Utility(\phi_b)
\right),
\]

where

\[
\phi_b
\sim
\Dir(\alpha_{\mathrm{root}}(b)).
\]

Estimate this by Monte Carlo. For \(m=1,\ldots,M\), sample

\[
\phi_b^{(m)}
\sim
\Dir(\alpha_{\mathrm{root}}(b)),
\qquad
b\in\Legal(s),
\]

then compute

\[
a_m^\star
=
\argmax_{b\in\Legal(s)}
\Utility(\phi_b^{(m)}).
\]

The estimator is

\[
\hat{\pi}_{\mathrm{search}}(a\mid s)
=
\frac{1}{M}
\sum_{m=1}^{M}
\one[a_m^\star=a].
\]

\begin{algorithmblock}{Algorithm 11: PosteriorBestPolicyTarget\((\alpha,M,\Legal(s))\)}

\Input action posteriors \(\alpha(a)\), number of samples \(M\), legal action set \(\Legal(s)\)

\Output Monte Carlo posterior-best policy target \(\hat{\pi}_{\mathrm{search}}\)

\begin{enumerate}[leftmargin=2em]
  \item Initialize counts:
  \[
  C(a)\gets 0
  \qquad
  \text{for all } a\in\Legal(s).
  \]

  \item For \(m=1,\ldots,M\):

  \begin{enumerate}[leftmargin=2em]
    \item For each legal action \(b\), sample
    \[
    \phi_b^{(m)}
    \sim
    \Dir(\alpha(b)).
    \]

    \item Compute
    \[
    a_m^\star
    \gets
    \argmax_{b\in\Legal(s)}
    \Utility(\phi_b^{(m)}).
    \]

    \item Increment
    \[
    C(a_m^\star)
    \gets
    C(a_m^\star)+1.
    \]
  \end{enumerate}

  \item Return
  \[
  \hat{\pi}_{\mathrm{search}}(a\mid s)
  =
  \frac{C(a)}{M}.
  \]
\end{enumerate}

\end{algorithmblock}

\section*{15. Final action selection}

Search produces a posterior-best policy target, but the committed action can be
selected in several ways.

First define the posterior mean scalar Q estimate:

\[
\hat{q}_T(a)
=
\Utility
\left(
\mean
\left(
\alpha_{\mathrm{root}}(a)
\right)
\right).
\]

Equivalently,

\[
\hat{q}_T(a)
=
\frac{
\alpha_{\mathrm{root},W}(a)
-
\alpha_{\mathrm{root},L}(a)
}{
\sum_z \alpha_{\mathrm{root},z}(a)
}.
\]

The supported committed-action modes are:

\[
\texttt{posterior\_sample},
\qquad
\texttt{posterior\_argmax},
\qquad
\texttt{scalar\_q\_argmax}.
\]

Posterior sampling chooses

\[
a_{\mathrm{play}}
\sim
\Cat
\left(
\hat{\pi}_{\mathrm{search}}(\cdot\mid s)
\right).
\]

Posterior argmax chooses

\[
a_{\mathrm{play}}
=
\argmax_a
\hat{\pi}_{\mathrm{search}}(a\mid s).
\]

Scalar-Q argmax chooses

\[
a_{\mathrm{play}}
=
\argmax_a
\hat{q}_T(a)
=
\argmax_a
\left(
\bar\phi_{a,W}-\bar\phi_{a,L}
\right),
\qquad
\bar\phi_a=\mean(\alpha_{\mathrm{root}}(a)).
\]

The posterior-best policy target is used for policy training regardless of which
committed-action mode is used.

\begin{algorithmblock}{Algorithm 12: FinishSearch\((\Tree,M,\texttt{commit})\)}

\Input completed posterior tree \(\Tree\), posterior-best samples \(M\), committed-action mode \(\texttt{commit}\)

\Output root posteriors, policy target, committed action

\begin{enumerate}[leftmargin=2em]
  \item For each legal root action \(a\), compute
  \[
  \alpha_{\mathrm{root}}(a)
  \gets
  \EdgePosterior(\Tree,v_{\mathrm{root}},a).
  \]

  \item Compute the posterior-best policy target:
  \[
  \hat{\pi}_{\mathrm{search}}
  \gets
  \PosteriorBestPolicyTarget
  \left(
  \alpha_{\mathrm{root}},
  M,
  \Legal(s_{\mathrm{root}})
  \right).
  \]

  \item If
  \[
  \texttt{commit}=\texttt{posterior\_sample},
  \]
  sample
  \[
  a_{\mathrm{play}}
  \sim
  \Cat
  \left(
  \hat{\pi}_{\mathrm{search}}
  \right).
  \]

  \item Else if
  \[
  \texttt{commit}=\texttt{posterior\_argmax},
  \]
  choose
  \[
  a_{\mathrm{play}}
  \gets
  \argmax_a
  \hat{\pi}_{\mathrm{search}}(a\mid s_{\mathrm{root}}).
  \]

  \item Else if
  \[
  \texttt{commit}=\texttt{scalar\_q\_argmax},
  \]
  choose
  \[
  a_{\mathrm{play}}
  \gets
  \argmax_a
  \Utility
  \left(
  \mean
  \left(
  \alpha_{\mathrm{root}}(a)
  \right)
  \right).
  \]

  \item \Return
  \[
  \left(
  \alpha_{\mathrm{root}},
  \hat{\pi}_{\mathrm{search}},
  a_{\mathrm{play}}
  \right).
  \]
\end{enumerate}

\end{algorithmblock}

\section*{16. Estimated-Q improvement view}

Let \(A_0\) be a Thompson-proposed root action from the initial posterior and
let \(C\) be a candidate set generated by Thompson search such that

\[
A_0 \in C.
\]

After search, define

\[
\hat{q}_T(a)
=
\Utility
\left(
\mean
\left(
\alpha_{\mathrm{root}}(a)
\right)
\right).
\]

If the committed action is

\[
A_+
=
\argmax_{a\in C}
\hat{q}_T(a),
\]

then the estimated-Q improvement is pointwise:

\[
\hat{q}_T(A_+)
\ge
\hat{q}_T(A_0).
\]

Taking expectation over Thompson proposal randomness gives

\[
\E
\left[
\hat{q}_T(A_+)
\right]
\ge
\E_{A_0}
\left[
\hat{q}_T(A_0)
\right].
\]

This is an estimated-value improvement statement. It becomes a true value
improvement only to the extent that \(\hat{q}_T\) is an accurate estimate of the
true action value.

\section*{17. Training targets}

For each searched root state \(s\), the stored training record should contain:

\[
s,
\qquad
\hat{\pi}_{\mathrm{search}}(\cdot\mid s),
\qquad
\alpha_{\mathrm{root}}(a)
\text{ for explored root actions},
\qquad
a_{\mathrm{play}},
\qquad
z \text{ when eventually known}.
\]

\subsection*{17.1 Policy target}

The policy target is

\[
\hat{\pi}_{\mathrm{search}}(\cdot\mid s).
\]

The policy loss is

\[
\mathcal{L}_\pi(s)
=
-
\sum_{a\in\Legal(s)}
\stopgrad
\left(
\hat{\pi}_{\mathrm{search}}(a\mid s)
\right)
\log
\pi_\theta(a\mid s).
\]

Equivalently, up to an additive constant independent of \(\theta\),

\[
\mathcal{L}_\pi(s)
=
D_{\mathrm{KL}}
\left(
\stopgrad
\left(
\hat{\pi}_{\mathrm{search}}(\cdot\mid s)
\right)
\,\|\,
\pi_\theta(\cdot\mid s)
\right).
\]

\subsection*{17.2 Q targets}

For an explored root action \(a\), the Q target is the completed root posterior:

\[
\beta_Q(s,a)
=
\alpha_{\mathrm{root}}(a).
\]

The Q loss is applied only to actions with positive completed evidence mass:

\[
\sum_z E_{v_{\mathrm{root}},a,z} > 0.
\]

Unexplored legal actions are masked out of the Q KL loss. This prevents the
network from being trained to reproduce its own unchanged prior as if it were a
search-improved target.

The Q loss is

\[
\mathcal{L}_Q(s)
=
\frac{
\sum_{a\in\Legal(s)}
m_Q(s,a)
D_{\mathrm{KL}}
\left(
\Dir
\left(
\stopgrad(\beta_Q(s,a))
\right)
\,\|\,
\Dir
\left(
\alpha_\theta^Q(s,a)
\right)
\right)
}{
\sum_{a\in\Legal(s)} m_Q(s,a)
+
\varepsilon
},
\]

where

\[
m_Q(s,a)
=
\one
\left[
\sum_z E_{v_{\mathrm{root}},a,z} > 0
\right].
\]

\subsection*{17.3 Value targets}

The value target may be constructed from eventual game outcome evidence or from
search/bootstrap evidence.

A search-based root value posterior proxy is

\[
\beta_{\mathrm{search}}^V(s)
=
\sum_{a\in\Legal(s)}
\hat\pi_{\mathrm{search}}(a\mid s)
\alpha_{\mathrm{root}}(a).
\]

This is the root-level analogue of the state posterior used during backup. It
uses \(E+\alpha^{\mathrm{base}}\) for each action, not only the raw evidence
\(E\).

If this proxy is used as bootstrap evidence, set

\[
d_V(s)
=
\mean
\left(
\beta_{\mathrm{search}}^V(s)
\right),
\qquad
\lambda_V
=
 c_V
\sum_z \beta_{\mathrm{search},z}^V(s),
\]

where \(c_V>0\) is a calibration coefficient. Then the posterior value target is

\[
\beta_V(s)
=
\alpha_{\theta^-}^V(s)
+
\lambda_V d_V(s).
\]

For a terminal outcome \(z\), use

\[
d_V(s)=e_z.
\]

The value loss is

\[
\mathcal{L}_V(s)
=
D_{\mathrm{KL}}
\left(
\Dir
\left(
\stopgrad(\beta_V(s))
\right)
\,\|\,
\Dir
\left(
\alpha_\theta^V(s)
\right)
\right).
\]

\begin{algorithmblock}{Algorithm 13: MakeTargets\((\Tree,\theta^-,M,\texttt{commit})\)}

\Input completed posterior tree \(\Tree\), data-generation parameters \(\theta^-\), posterior-best samples \(M\), committed-action mode \(\texttt{commit}\)

\Output training targets for the root state

\begin{enumerate}[leftmargin=2em]
  \item Run
  \[
  \left(
  \alpha_{\mathrm{root}},
  \hat{\pi}_{\mathrm{search}},
  a_{\mathrm{play}}
  \right)
  \gets
  \FinishSearch(\Tree,M,\texttt{commit}).
  \]

  \item Store policy target:
  \[
  \pi_{\mathrm{target}}(\cdot\mid s)
  \gets
  \hat{\pi}_{\mathrm{search}}(\cdot\mid s).
  \]

  \item For each legal root action \(a\), define Q mask
  \[
  m_Q(s,a)
  \gets
  \one
  \left[
  \sum_z E_{v_{\mathrm{root}},a,z} > 0
  \right].
  \]

  \item For each explored root action, store
  \[
  \beta_Q(s,a)
  \gets
  \alpha_{\mathrm{root}}(a).
  \]

  \item Optionally store the search-weighted value posterior proxy:
  \[
  \beta_{\mathrm{search}}^V(s)
  \gets
  \sum_{a\in\Legal(s)}
  \hat\pi_{\mathrm{search}}(a\mid s)
  \alpha_{\mathrm{root}}(a).
  \]

  \item Store
  \[
  a_{\mathrm{play}}
  \]
  for self-play.

  \item Store enough information to construct the eventual value target once
  outcome or bootstrap evidence is available.

  \item \Return training targets.
\end{enumerate}

\end{algorithmblock}

\subsection*{17.4 Flattening searched tree nodes}

The root trajectory record above is not the only possible source of training
signal. After search has completed for a self-play root, the completed posterior
tree may also be flattened into additional per-state training records.

For a tree node \(v\), let \(s_v\) be the represented game state. A flattened
tree-node record has the same shape as a root record:

\[
s_v,
\qquad
\hat{\pi}_{\mathrm{search}}(\cdot\mid s_v),
\qquad
\beta_Q(s_v,a),
\qquad
\beta_V(s_v),
\qquad
m_\pi(s_v),
\qquad
m_V(s_v),
\qquad
m_{\mathrm{out}}(s_v).
\]

For an expanded non-terminal node with completed child evidence, define

\[
\beta_Q(s_v,a)
=
\alpha^{\mathrm{post}}_{v,a}
=
\alpha^{\mathrm{base}}_{v,a}
+
E_{v,a}.
\]

The node is eligible for policy and Q training only if it has positive completed
evidence:

\[
\sum_{a\in\Legal(s_v)}
\sum_z E_{v,a,z}
>
\epsilon_{\mathrm{tree}}.
\]

This excludes non-terminal leaves that were merely evaluated by the network but
never received completed child evidence. Such leaves mostly reproduce the
network's own prior and are therefore weak self-distillation rather than
search-improved supervision.

For eligible expanded nodes, the local policy target is computed in the same way
as at the root:

\[
\hat{\pi}_{\mathrm{search}}(\cdot\mid s_v)
=
\PosteriorBestPolicyTarget
\left(
\left\{
\alpha^{\mathrm{post}}_{v,a}
\right\}_{a\in\Legal(s_v)},
M
\right).
\]

The value posterior proxy for the same node is

\[
\beta_V(s_v)
=
\alpha_{\theta^-}^V(s_v)
+
c_V
\sum_{a\in\Legal(s_v)}
\hat{\pi}_{\mathrm{search}}(a\mid s_v)
\beta_Q(s_v,a).
\]

These internal expanded-node rows use

\[
m_\pi(s_v)=1,
\qquad
m_V(s_v)=1,
\qquad
m_{\mathrm{out}}(s_v)=0,
\]

so that policy, Q-Dirichlet, and value-Dirichlet losses apply, but hard
eventual-outcome losses do not. They do not lie on the played trajectory, so a
final game result is not generally available for them.

Terminal tree leaves may also be stored. For a terminal leaf with outcome \(z\),
use a value-only record:

\[
\beta_V(s_v)=e_z,
\qquad
m_\pi(s_v)=0,
\qquad
m_V(s_v)=1,
\qquad
m_{\mathrm{out}}(s_v)=1.
\]

The root node itself is usually not duplicated in the flattened tree records,
because the played self-play trajectory already stores the root record. It may be
included explicitly if a particular implementation wants to make the flattened
tree batch self-contained.

In a batched implementation, each self-play move can export a fixed-capacity
tree-node array. Unused slots are padded with zero data and false masks:

\[
m_\pi=m_V=m_{\mathrm{out}}=0.
\]

After this export, the tree structure is discarded for training purposes. The
learner sees only a flat collection of rows with observations, targets, legal
masks, and loss masks.

\section*{18. Practical implementation notes}

\subsection*{18.1 CPU tree layout}

The CPU tree should be implemented as an arena of nodes and edges. Prefer
integer indices over pointers:

\[
\texttt{node\_id},
\qquad
\texttt{edge\_id},
\qquad
\texttt{child\_id}.
\]

For small or medium action spaces such as Hex, dense edge arrays are often
reasonable:

\[
\texttt{nodes} \times |\Actions| \times 3.
\]

For large action spaces such as chess-style encodings, sparse edge storage may
be required.

\subsection*{18.2 Redis and tree storage}

Redis can be useful around the tree, but it should not be the inner-loop tree
store for selection and backup.

Use an in-process arena for active trees because search repeatedly reads and
writes adjacent node and edge statistics. Redis round trips are too expensive
for per-edge Thompson selection, posterior updates, and bottom-up backup.

Reasonable Redis uses are:

\begin{enumerate}[leftmargin=2em]
  \item GPU work queues and result queues.
  \item Self-play trajectory queues.
  \item Root summaries after search.
  \item Checkpoints or serialized trees when latency is not critical.
  \item Optional transposition or opening caches, if treated as approximate and
  not queried on every edge update.
\end{enumerate}

The practical split is:

\[
\text{active tree in RAM},
\qquad
\text{coordination and persistence in Redis}.
\]

\subsection*{18.3 GPU boundary}

The GPU evaluator should be the compiled part:

\[
\GpuEvaluate(\mathcal{Q},\theta).
\]

It receives a preallocated batch of observations and legal masks, and returns
preallocated arrays containing:

\[
\ell_\theta(s,\cdot),
\qquad
\alpha_\theta^V(s),
\qquad
\alpha_\theta^Q(s,\cdot).
\]

The CPU tree should not wait after every single descent. It should fill batches
using both game-parallelism and leaf-parallelism.

\subsection*{18.4 Flattened training batch boundary}

The mutable posterior tree is a search-side data structure. It is not sent to the
training step as a tree. Instead, after a search finishes, the CPU or host-side
search code materializes dense training arrays:

\[
\texttt{obs},
\qquad
\texttt{policy\_target},
\qquad
\texttt{legal\_mask},
\qquad
\texttt{beta\_Q},
\qquad
\texttt{beta\_V},
\qquad
\texttt{loss\_masks}.
\]

The root trajectory rows and flattened tree rows are then concatenated, shuffled,
trimmed to a multiple of the training batch size, and reshaped into minibatches.
At that point there is no parent/child structure left in the learner input.

For a self-play batch size \(B\), game horizon \(T\), and per-root simulation
count \(N\), the root-only data volume is

\[
B\,T.
\]

If all tree nodes were retained, the rough upper bound would be

\[
B\,T\,(N+1),
\]

because each completed simulation can add at most one new node in addition to the
root. In practice, filtering non-terminal leaves without completed evidence and
deduplicating transpositions reduces this count.

\subsection*{18.5 Tree lifetime and reuse}

In the arena implementation, each posterior tree is temporary. For each self-play
move, search builds a fresh arena from the current batch of root states, runs the
configured number of simulations, extracts root targets and optional flattened
tree-node rows, commits one action per root, and then drops the arena.

Thus the active tree is not reused across self-play moves or training
iterations. Only the neural network parameters, optimizer state, and checkpointed
training state persist. A separate store-backed implementation may cache
serialized nodes or root summaries, but that is a search-system optimization, not
part of the basic training-data contract.

\subsection*{18.6 Recommended initial knobs}

A reasonable initial setting is

\[
L \in \{4,8\},
\]

where \(L\) is the number of in-flight leaves allowed per root tree.

For the search-weighted state backup, start conservatively:

\[
c_{\mathrm{state}} \in [0.01,0.25].
\]

If the backup is too weak, increase \(c_{\mathrm{state}}\). If the tree becomes
overconfident or collapses too quickly to one line, decrease it.

The terminal evidence strength should usually be larger than the neural leaf
evidence strength:

\[
c_{\mathrm{terminal}} > c_{\mathrm{leaf}}.
\]

\subsection*{18.7 Completed evidence only}

Only completed evidence \(E_{v,a}\) is used in posteriors and training targets:

\[
\alpha^{\mathrm{post}}_{v,a}
=
\alpha^{\mathrm{base}}_{v,a}
+
E_{v,a}.
\]

In-flight scheduling state is not Bayesian evidence and is not included in
\(\alpha^{\mathrm{post}}_{v,a}\).

\end{document}
